\documentclass[11pt]{article}
\usepackage[final]{acl}
\usepackage{times}
\usepackage{latexsym}
\usepackage[utf8]{inputenc}
\usepackage{microtype}
\usepackage{amsmath,amssymb}
\usepackage{booktabs}
\usepackage{graphicx}
\usepackage{url}

\title{Đề tài 08: Nghiên cứu và Xây dựng Mô hình Dịch máy Anh--Việt \\ Dựa trên Kiến trúc Seq2Seq Attention và Transformer}

\author{
  Lê Minh Hiếu \\
  Trường Đại học Bách khoa Hà Nội \\
  \texttt{hieu.lm@sis.hust.edu.vn}
}

\date{}

\begin{document}
\maketitle

\begin{abstract}
Dịch máy nơ-ron (Neural Machine Translation -- NMT) cho cặp ngôn ngữ Anh--Việt đặt ra nhiều thách thức do sự dị biệt sâu sắc về ngữ pháp, trật tự từ và đặc tính đơn lập của tiếng Việt. Đề xuất này trình bày kế hoạch nghiên cứu và tự cài đặt từ đầu (from scratch) hai kiến trúc dịch máy cốt lõi: mô hình Seq2Seq LSTM tích hợp Attention và mô hình Transformer nhỏ. Toàn bộ thực nghiệm được kiểm soát nghiêm ngặt trên bộ dữ liệu IWSLT'15 En--Vi và đánh giá chuẩn hóa bằng chỉ số SacreBLEU.
\end{abstract}

\section{Mô tả Bài toán (Problem Formulation)}
Bài toán đặt ra là ánh xạ một chuỗi token tiếng Anh nguồn $X = (x_1, x_2, \dots, x_S)$ sang chuỗi token tiếng Việt đích tương ứng $Y = (y_1, y_2, \dots, y_T)$.
Mục tiêu là tối ưu hóa phân phối xác suất có điều kiện:
\begin{equation}
\hat{Y} = \arg\max_{Y} P(Y \mid X) = \arg\max_{Y} \prod_{t=1}^{T} P(y_t \mid y_{<t}, X; \theta)
\end{equation}
Bên cạnh chất lượng dịch thuật, mô hình yêu cầu trích xuất ma trận trọng số chú ý (Attention Alignment Matrix) $A \in \mathbb{R}^{T \times S}$ nhằm phục vụ việc phân tích gióng hàng ngữ nghĩa và giải thích lỗi dịch (Error Analysis).

\section{Dữ liệu & Tiền xử lý (Data Pipeline)}
Dự án sử dụng bộ dữ liệu chuẩn quốc tế \textbf{IWSLT'15 English--Vietnamese} \cite{luong2015effective} gồm $\approx 133,317$ cặp câu song ngữ từ các bài diễn thuyết TED Talks.
\begin{itemize}
    \item \textbf{Phân chia tập dữ liệu:} Tập Train (132,418 câu sau khi lọc), tập Validation (\texttt{tst2012}, 1,550 câu) và tập Test (\texttt{tst2013}, 1,262 câu).
    \item \textbf{Tiền xử lý:} Chuẩn hóa Unicode NFC, chuyển chữ thường (lowercasing), lọc bỏ ký tự đặc biệt và giới hạn độ dài câu trong khoảng $[1, 100]$ từ để kiểm soát bộ nhớ GPU.
    \item \textbf{Bộ tách từ Subword:} Áp dụng thuật toán Byte Pair Encoding (BPE) \cite{sennrich2016neural} dùng chung (Joint Vocabulary) cho cả hai ngôn ngữ với kích thước từ vựng $V = 8,000$ tokens, hạn chế triệt để hiện tượng từ ngoài từ điển (Out-Of-Vocabulary -- OOV).
\end{itemize}

\section{Mô hình Dự kiến (Proposed Architectures)}
\paragraph{1. Baseline -- Seq2Seq LSTM có Attention:}
Tự cài đặt mô hình Encoder--Decoder dựa trên mạng nơ-ron hồi quy LSTM \cite{sutskever2014sequence} kết hợp cơ chế chú ý Bahdanau \cite{bahdanau2014neural} hoặc Luong \cite{luong2015effective}. Cơ chế attention cho phép bộ giải mã trích xuất vector ngữ cảnh động $c_t = \sum_{s=1}^{S} \alpha_{ts} h_s$ tại mỗi bước sinh từ thay vì phụ thuộc vào một vector ngữ cảnh tĩnh.

\paragraph{2. Mô hình Chính -- Transformer nhỏ (From Scratch):}
Cài đặt hoàn chỉnh kiến trúc Transformer theo \citet{vaswani2017attention} bao gồm: Positional Encoding dạng sin/cos, Multi-Head Scaled Dot-Product Attention, Position-wise Feed-Forward Network, Residual Connection và Layer Normalization. Cài đặt chuẩn xác Padding Mask và Causal Mask (Look-ahead mask) để ngăn chặn rò rỉ thông tin tương lai ở Decoder.

\section{Phương pháp Đánh giá (Evaluation Metrics)}
Chất lượng dịch máy được lượng hóa tự động bằng chỉ số \textbf{BLEU} chuẩn quốc tế thông qua thư viện \texttt{sacrebleu} \cite{post2018call}. Quá trình suy luận sẽ tiến hành so sánh đối sánh giữa hai chiến lược giải mã: \textit{Greedy Search} và \textit{Beam Search} (kích thước beam $k \in \{3, 5\}$). Điểm BLEU kỳ vọng đạt từ 22--26 điểm cho Baseline và 27--32 điểm cho Transformer.

\section{Kế hoạch Triển khai (Work Breakdown)}
Dự án được triển khai độc lập trên nhánh \texttt{Hiếu} qua 6 tuần:
\begin{itemize}
    \item \textbf{Tuần 1 (Đã hoàn thành):} Thiết lập môi trường, xây dựng pipeline dữ liệu (Preprocess, BPE Tokenizer, DataLoader dynamic padding), chuẩn hóa metric BLEU và soạn thảo Proposal.
    \item \textbf{Tuần 2:} Tự cài đặt và huấn luyện Baseline Seq2Seq LSTM có Attention; tích hợp log W\&B/TensorBoard.
    \item \textbf{Tuần 3:} Tự cài đặt Transformer nhỏ từ đầu; huấn luyện và nộp Milestone Report (tối đa 4 trang).
    \item \textbf{Tuần 4:} Thí nghiệm đối sánh có kiểm soát và thực hiện Ablation Study (số layer, số head, dropout).
    \item \textbf{Tuần 5:} Trực quan hóa Attention Map, phân tích định tính lỗi sai và hoàn thiện bản nháp báo cáo.
    \item \textbf{Tuần 6:} Hoàn thiện báo cáo cuối (8 trang chuẩn ACL), đóng gói mã nguồn, làm demo và chuẩn bị vấn đáp.
\end{itemize}

\bibliography{references}

\end{document}
